\subsection{One retained vector for the two short-arm sums}
\label{sec:unit-retained-square}
Use the original finite kernel, with $B=-K$, $u=B\one$, and
$\mathsf X=-B^2$, and write $g=K\one=-u$. The path coefficients $\beta_j$ are those in
\eqref{eq:Q188:conv}; put $\tau=d-\delta_2$, $\theta=13/32$,
and $\kappa=1-\theta$.

\begin{lemma}[A retained pair of arm sums]
\label{lem:unit-retained-pair}
Let $n\ge5$. Take either $\sigma=2$, $0\le d\le1/3$, or
$\sigma=9/4$, $0\le d\le1/5$, and set $m=1-\sigma d/2$.
For $n=5$ put
\[
 a_1=d,\qquad a_2=(\sigma-2)d+\delta_2,\qquad a_3=0.
\]
For $n\ge6$ put
\[
 \lambda=\frac{\sigma d-\beta_n-\beta_{n-1}}
                    {\beta_{n-2}+(1-d)\beta_{n-3}},\qquad
 a_1=\beta_{n-1},\quad a_2=\lambda\beta_{n-2},\quad
 a_3=\lambda\beta_{n-3}.
\]
At $d=0$ all $a_i$ are zero. These are genuine retained coefficients:
$0\le a_i\le\beta_{n-i}$. Define
\[
 C_0=a_1+a_2+a_3,\quad D_0=a_2+a_3,\quad E_0=a_3,\qquad
 U=\sum_{j=3}^{n}K^j\one-C_0g-D_0Kg-E_0K^2g.
\]
The same actual vector gives both comparisons
\begin{equation}\label{eq:unit-retained:dual-sums}
 h_2+f_n\le2m+U,\qquad f_2+h_n\le2m-U.
\end{equation}
Each right side is nonnegative. In particular
$h_2f_n\le(m+U/2)^2$ and $f_2h_n\le(m-U/2)^2$.

Write $U=[a(\mathsf X)-b(\mathsf X)B]u$. At $n=5$ its exact
coordinates are
\[
 a=(\sigma-1)d+\delta_2+x-x^2,\qquad
 b=(\sigma-2)d+\delta_2+x.
\]
For $n\ge6$, writing $G_{n-2}(i\sqrt x)=a_N+i\sqrt x\,b_N$,
they are
\begin{equation}\label{eq:unit-retained:coordinates}
 a=C_0+x(a_N-E_0),\qquad b=D_0+xb_N,
 \qquad a-(1-x)b\ge a_1+a_2x\ge\tau.
\end{equation}
All statements retain the global $d,\delta_2,E_0$ while $x$ varies.
\end{lemma}
\begin{proof}
For $d>0$ the denominator defining $\lambda$ is at least
$(2-d)\tau>0$. Its numerator is nonnegative. The full available
constant loss is at least $(4-d)\tau$: it exceeds $2d$ on
$[0,1/3]$, since $209/96>2$, and exceeds $9d/4$ on $[0,1/5]$,
since $361/160>9/4$. Thus $0\le\lambda\le1$.
For $n=5$, $a_2\le(21/32)d\le\beta_3$ even at $\sigma=9/4$;
$a_3=0$ takes nothing from $\beta_2=0$.
In both cases
\[
 \beta_n+a_1+a_2+(1-d)a_3=\sigma d,
 \quad C_0=\sigma d-\beta_n+dE_0,
 \quad D_0=\sigma d-\beta_n-\beta_{n-1}+dE_0.
\]
Retain $\beta_n\one,a_1f_1,a_2f_2,a_3f_3$ in the positive
convolution for $f_n$. Since $h_2+\sum_{j=0}^nK^j\one
=2+\sum_{j=3}^nK^j\one$, this proves the first sum.
Applying $T$ to the convolution for $f_{n-1}$ gives exactly
\[
 f_2+h_n=2-\beta_n-\sum_{j=3}^nK^j\one
             -\sum_{j=1}^{n-3}\beta_{n-j}h_j.
\]
Retaining $a_1h_1,a_2h_2,a_3h_3$ proves the second sum, with
the opposite sign of the same $U$.
The finite geometric identity gives the displayed coordinates.
For $N\ge4$, alternating sums give
$1-x\le a_N,b_N\le1-x+x^2$ and $a_N\ge(1-x)b_N$.
The last inequality in \eqref{eq:unit-retained:coordinates} follows
after inserting $C_0-D_0=a_1$ and $D_0-E_0=a_2$;
the $n=5$ formula verifies the same assertion directly.
\end{proof}

\begin{lemma}[The two retained first-row packets]
\label{lem:reference-retained-packet}
Take the actual vector $U=[a(\mathsf X)-b(\mathsf X)B]u$ of
Lemma~\ref{lem:unit-retained-pair}, with $n=c\ge5$. The following
two rows hold for $0<d\le d_{\max}$:
\[
\begin{array}{c|c|c|c|c|c}
 \sigma&d_{\max}&\ell&a\text{ upper}&b\text{ upper}&D\\\hline
 2&1/4&2&7/10&3/5&d/2\\
 9/4&1/5&15/7&16/25&9/16&2d/5.
\end{array}
\]
In each row, with $m=1-\sigma d/2$, $k=m^{-1}$ and $R=-kU$, one has
\begin{equation}\label{eq:reference-retained:coefficient}
 \ell(a^2+xb^2)\le a+b,\qquad
 T(qf_2h_c)\le m^2(\bar f_2+BR+D\one).
\end{equation}
The same coefficients satisfy $a,b\ge x(1-x)$ and
$ka\ge k\tau+(1-x)kb$.
Let $D_i=\E B^iR$ and $\gamma_j=\beta_j-D_{j-2}$. Their exact signs and
the bounds needed at every later finite index are
\begin{equation}\label{eq:reference-retained:signed-moments}
\begin{aligned}
 D_{4j+1}&=k\int x^{2j}a\,d\nu,&
 D_{4j+2}&=k\int x^{2j+1}b\,d\nu,\\
 D_{4j+3}&=-k\int x^{2j+1}a\,d\nu,&
 D_{4j+4}&=-k\int x^{2j+2}b\,d\nu,\\
 D_1&<d,&D_2&<d/3,\\
 \gamma_3&\ge0,&\gamma_4&\ge2d/3,\\
 \gamma_5&\ge\tau,&\gamma_j&\ge41d/96\quad(j\ge5).
\end{aligned}
\end{equation}
Every positive $D_i$ with $i\ge3$ is less than $d/6$. In the $9/4$ row,
also $D_1+D_2<5d/4$. The $2d$ row retains its own original vector;
it is available through $d=1/4$ independently of the $9d/4$ row.
In the $9d/4$ row, the displayed coordinate cap further gives
$D_1<128d/155$, hence $\gamma_3>27d/155>d/8$.
At $d=0$, $U=R=0$ and all non-strict bounds extend directly.
\end{lemma}
\begin{proof}
Write $E=a_3$, $C=a_1+a_2+a_3$ and $v(x)=x-x^2$, $w(x)=x-x^2+x^3$.
The actual retained coefficients give
\[
 (\sigma-1)d+dE\le C\le(\sigma-\kappa)d+dE,
 \quad 0\le E\le d,\quad
 a\le C+w(x)-xE,\quad b\ge C-d+v(x)\ge0.
\]
The last two comparisons also hold at $c=5$, using its explicit
coordinates. For $c=5$, its explicit coordinates are increasing in $x$ and give
\[
 a\le\theta-\theta^2+(\sigma-1+\theta)d,\qquad
 b\le\theta+(\sigma-2+\theta)d.
\]
At the respective $d_{\max}$ these bounds are
$(607/1024,65/128)$ for $\sigma=2$, and
$(2931/5120,43/80)$ for $\sigma=9/4$; they are strictly below
the four coordinate caps in the table.
For $c\ge6$, first increase $x$ to $\theta$, keeping the original
$d,E$ fixed. This is legitimate since
$w'(x)-E\ge2/3-E>0$. Because $d<\theta$,
\[
 a\le w(\theta)+(\sigma-\kappa)d,\qquad
 b\le w(\theta)+(\sigma-2\kappa)d+d^2.
\]
At $d=d_{\max}$ the respective bounding pairs are
\[
\begin{aligned}
 \sigma=2 &: \quad(21621/32768,18805/32768),\\
 \sigma=9/4 &: \quad(104777/163840,459373/819200).
\end{aligned}
\]
They imply the displayed coordinate bounds and $0\le ka,kb<1$.
Also $b\ge x(1-x)$, while the lower bound for $a$ exceeds
$x(1-x)+d(\sigma-1)+E(d-x)\ge x(1-x)$.
The relation between $ka$ and $kb$ follows from
\eqref{eq:unit-retained:coordinates}.

Put $\Phi(a,b)=\ell(a^2+xb^2)-a-b$. The displayed bounds imply
$2\ell xb<1$, so $\Phi$ decreases in $b$ and is convex in $a$.
The endpoint $a=0$ is nonpositive. At the other endpoint put
$A=C+w(x)-xE$, and decrease $b$ to $B_*=C-d+v(x)$. For fixed
original $C,d,E$, the exact second derivative is
\[
 2\ell\bigl[(1-E-2x+3x^2)^2
 +(6x-2)(d-Ex+x^3)+x(1-2x)^2\bigr]+4-6x.
\]
For $x\le1/3$, the square is at least $25/144$, the $d$ term at least
$-1/2$, and the $-Ex$ term is nonnegative. The remaining terms combine as
$x(1-4x+2x^2+6x^3)\ge0$; its parenthesis decreases to $1/9$.
Thus this derivative is at least $101/168>0$ for both rows.
For $x\ge1/3$ the bracket is nonnegative. It suffices to check
$x=0,\theta$.

At zero, $\Phi=\ell C^2-2C+d$ decreases in $C$, since
$C\le53/128$ and $(15/7)(53/128)<1$.
Substituting $C\ge(\sigma-1)d$ proves the required sign.
At $x=\theta$, convexity reduces $C$ to
$\alpha d+dE$, where $\alpha=\sigma-1,\sigma-\kappa$, and then
$E$ to $0,d$. The $E=0$ expressions are convex in $d$.
For $E=d$ the coordinates are
\[
 A=w(\theta)+(\alpha-\theta)d+d^2,\qquad
 B_*=v(\theta)+(\alpha-1)d+d^2.
\]
Their second $d$ derivative is
$2\ell[(A')^2+2A+\theta((B_*')^2+2B_*)]-4>0$, because
$\kappa^2+2w(\theta)+2\theta v(\theta)>1$ and $\alpha\ge1$.
Only $d=0,d_{\max}$ remain. At zero use
$w(\theta)<5/16$, $v(\theta)<1/4$ and
$w(\theta)+v(\theta)>1/2$.
For $d=d_{\max}$, in the order
$(\alpha,E)=(\sigma-1,0),(\sigma-\kappa,0),
(\sigma-1,d),(\sigma-\kappa,d)$, exact substitution gives
\[
 \begin{aligned}
 -536870912\,\Phi\big|_{\sigma=2}
   &\in\{69198055,19544295,112116455,65444583\},\\
 -187904819200\,\Phi\big|_{\sigma=9/4}
   &\in\{20259598945,2787865185,33811712609,17975757409\}.
 \end{aligned}
\]
This proves the coefficient inequality on the full continuous domains.

Let $V=\E U^2=\int(a^2+xb^2)\,d\nu$.
Expanding the original positive square gives
\[
 \frac{T(qf_2h_c)}{m^2}
 \le\bar f_2+BR+\delta\one,\qquad
 \delta=-k\int(a+b)\,d\nu+k\sqrt{dV}+k^2V.
\]
Here the linear row error is bounded by $\E|uR|$, and the quadratic
term by $\E R^2$, using the original kernel row bound and $q\le2$.
For $\sigma=2$, use
$\sqrt{dV}\le3d/8+(2/3)V$ and $k+2/3\le2$ to get
$\delta\le3kd/8\le d/2$.
For $\sigma=9/4$, use
$\sqrt{dV}\le3d/10+(5/6)V$ and $k+5/6\le15/7$ to get
$\delta\le3kd/10\le12d/31<2d/5$.

Skew adjunction gives the four exact sign classes in
\eqref{eq:reference-retained:signed-moments}. The coordinate caps in
the table imply $D_1<d$, $D_2<d/3$ and, in the second row,
$D_1+D_2<5d/4$. The later positive classes cost at most
$k\theta^2(\sup a)d$ and $k\theta^3(\sup b)d$, both less than $d/6$.
Combining them with $\beta_j\ge\tau\ge\kappa d$ proves all the
$\gamma$ bounds. At $d=0$ the retained vectors vanish.
\end{proof}

\begin{lemma}[Positive continuation of a retained square]
\label{lem:reference-retained:continuation}
For either row above, let $F=qf_2h_c$ and use its same actual $R,D$.
For every finite $L\ge1$,
\begin{equation}\label{eq:reference-retained:continuation}
\begin{aligned}
 T^LF&\le m^2Y_L,\\
 Y_L&=\bar f_{L+1}+B^LR+D\bar f_{L-1}
             +\sum_{i=1}^{L-1}D_i\bar f_{L-1-i}\\
 &=\one+G_{L+1}(B)u+B^LR+D\bar f_{L-1}
             -\sum_{j=3}^{L+1}\gamma_j\bar f_{L+1-j}.
\end{aligned}
\end{equation}
Every loss in the last sum is nonnegative.
\end{lemma}
\begin{proof}
Apply $T^{L-1}$ to the original first row and use $T=J+B$.
The identity
$T^{L-1}(BR)=B^LR+\sum_{i=1}^{L-1}(\E B^iR)\bar f_{L-1-i}$
follows by finite induction. Substituting the positive path convolution
for $\bar f_{L+1}$ gives the second expression.
\end{proof}

\begin{lemma}[Strong norms for the $9d/4$ packet]
\label{lem:reference-retained:strong-norm}
For the actual $\sigma=9/4$ coefficients of
Lemma~\ref{lem:unit-retained-pair}, for every finite $c\ge5$ and
$0\le d\le1/5$, one has
\begin{equation}\label{eq:reference-retained:strong-norm}
 \sqrt{a^2+xb^2}\le u_0=3/8+9d/5\le147/200,
 \qquad b\le z_0=\theta+13d/16\le91/160.
\end{equation}
\end{lemma}
\begin{proof}
For $c=5$ the coordinates increase in $x$. First set $x=\theta$
with the original $d,\delta_2$ fixed, then use $\delta_2\le\theta d$.
The resulting vector is bounded coordinatewise by
\[
 (\theta-\theta^2,\sqrt\theta\,\theta)
   +d(53/32,\sqrt\theta\,21/32).
\]
Its constant and linear squared norms are
$131313/1048576<9/64$ and $95621/32768<81/25$.
The triangle inequality proves the norm bound;
$b\le\theta+21d/32\le z_0$ gives the second assertion.

For $c\ge6$ put $w(x)=x-x^2+x^3$. The actual global moment gives
\[
 a\le w(x)+5d/4+\delta_2+E(d-x),\qquad
 b\le w(x)+d/4+2\delta_2+dE.
\]
Both upper functions increase in $x$: their derivatives are at least
$2/3-E\ge7/15$ and $2/3$. Set $x=\theta$, then
$\delta_2\le\theta d$. With $w_\theta=10101/32768<5/16$,
the bounding coordinates are
\[
 a_*=w_\theta+(53/32)d+E(d-\theta),\qquad
 b_*=w_\theta+(17/16)d+dE.
\]
Their squared norm is convex in $E\in[0,d]$.
At $E=0$ the constant norm is below $3/8$ and the linear
squared norm is $26229/8192<81/25$.
At $E=d$ the linear squared norm is $16557/8192<100/49$ and
the quadratic vector has norm $\sqrt{45/32}<6/5$. Thus
\[
 \sqrt{a_*^2+\theta b_*^2}
 \le3/8+(10/7)d+(6/5)d^2
 \le3/8+(10/7+6/25)d<3/8+9d/5.
\]
Finally $b\le5/16+(17/16)d+d^2$, whose excess above $z_0$ is
$-3/32+d/4+d^2\le-3/800$. All frequency maximizations kept the
original $d,\delta_2,E$ fixed.
\end{proof}

\begin{lemma}[The two short retained polynomials]
\label{lem:reference-retained:short-norms}
For the same $\sigma=9/4$ packet, put $t=ka$, $s=kb$ and
\[
 Q_2=1-sx+z(1-t),\qquad
 Q_3=1-x+tx+z(1-sx).
\]
For every actual finite $c\ge5$ and every $x\in[0,\theta]$,
\begin{equation}\label{eq:retained-reference:Q2-tight}
 |Q_2|^2\le1+8/63<8/7,
\end{equation}
\begin{equation}\label{eq:retained-reference:Q3-tight}
 |Q_3|^2\le8/7.
\end{equation}
\end{lemma}
\begin{proof}
Both normalized imaginary coordinates are nonnegative since
$0\le t,s\le1$. Using $a,b\ge x(1-x)$ and $k\ge1$ gives
\[
 |Q_2|^2\le1+x-4x^2+5x^3-x^4-x^5+x^6
 \le1+x-(63/32)x^2\le1+8/63.
\]
Indeed $4-5x+x^2+x^3-x^4$ decreases to
$2279183/1048576>63/32$.

For $Q_3$, first suppose $x\le d$.
Then $a\le(53/32)d+dE+x(1-E)\le85d/32$, hence
$ka\le85/124<3/4$. Thus
$|Q_3|^2\le(1-x/4)^2+x\le441/400<8/7$.
For $x\ge d$, increase $a$ to $C+w(x)-xE$ and decrease
$b$ to $C-d+v(x)$, with $v=x-x^2$, $w=x-x^2+x^3$.
This increases both nonnegative coordinates. Their squared norm
increases in $C$ for $x>0$, because its derivative is
$2kx[\operatorname{Re}Q_3-x\operatorname{Im}Q_3/\sqrt x]>0$;
the parenthesis is at least $1-2x>0$.
The case $x=0$ is direct.
Take $C=53d/32+dE$. Increasing $E$ now decreases both coordinates,
so take $E=0$. At fixed $x$ the coordinates are affine in
$d/m\in[0,8/31]$, and the squared norm is convex. It suffices to
use $d=0,1/5$, allowing the whole rectangle for this last bound.
At zero it is at most $(1-x+x^2)^2+x\le1$.
At $1/5$, the coordinates are
\[
 1-(71/124)x+(40/31)x^2-(40/31)x^3+(40/31)x^4,
 \quad 1-(21/124)x-(40/31)x^2+(40/31)x^3.
\]
In their squared norm $P(x)=1+\sum_{j=1}^8 a_jx^j$,
\[
\begin{aligned}
 (a_1,\ldots,a_5)
 &=(-9/62,39513/15376,-101639/15376,\\
 &\hspace{2em}134400/15376,-55040/15376),\\
 (a_6,a_7,a_8)&=(k^2,-k^2,k^2),\qquad k=40/31.
\end{aligned}
\]
Thus $P''\ge5-40x+104x^2-72x^3>0$; the remaining term is
$k^2x^4(30-42x+56x^2)>0$.
For $y=x-1/4\le5/32$, the cubic lower bound is
$3/8-(3/2)y+50y^2-72y^3
\ge3/8-(3/2)y+(155/4)y^2>0$.
At $x=0$ the norm is one; at $\theta$ the two coordinates are
less than $93/100,81/100$, giving
$P(\theta)<362061/320000<8/7$.
Convexity proves the bound on the entire frequency interval.
\end{proof}
